
\documentclass{article}
\usepackage{colortbl}
\usepackage{makecell}
\usepackage{multirow}
\usepackage{supertabular}

\begin{document}

\newcounter{utterance}

\centering \large Interaction Transcript for game `cladder', experiment `full\_v1.5\_default', episode 3923 with selfrefine{-}llama{-}3.1{-}8b{-}instruct.
\vspace{24pt}

{ \footnotesize  \setcounter{utterance}{1}
\setlength{\tabcolsep}{0pt}
\begin{supertabular}{c@{$\;$}|p{.15\linewidth}@{}p{.15\linewidth}p{.15\linewidth}p{.15\linewidth}p{.15\linewidth}p{.15\linewidth}}
    \# & \multicolumn{2}{c}{Player} && \multicolumn{2}{c}{Game Master} \\
    \hline

    \theutterance \stepcounter{utterance}  
    & & & \multicolumn{4}{p{0.6\linewidth}}{
        \cellcolor[rgb]{0.9,0.9,0.9}{
            \makecell[{{p{\linewidth}}}]{
                \texttt{\tiny{[P1$\langle$GM]}}
                \texttt{You are an expert in causal inference. The following question is not a typical commonsense query, but rather a meticulously designed question created by a professor specializing in causal inference, intended to assess the students' mastery of the course content.} \\
\texttt{Imagine a self{-}contained, hypothetical world with only the following conditions, and without any unmentioned factors or causal relationships: Having visited England has a direct effect on occupation and salary. Skill has a direct effect on occupation and salary. Occupation has a direct effect on salary. Skill is unobserved. The overall probability of having visited England is 43\%. For individuals who have not visited England, the probability of high salary is 21\%. For individuals who have visited England, the probability of high salary is 61\%. Is high salary less likely than low salary overall?} \\
\texttt{Start your answer with 'Yes' or 'No', followed by additional reasoning or evidence to support your explanation.} \\
            }
        }
    }
     \\ \\

    \theutterance \stepcounter{utterance}  
    & & \multicolumn{4}{p{0.6\linewidth}}{
        \cellcolor[rgb]{0.95,0.95,0.95}{
            \makecell[{{p{\linewidth}}}]{
                \texttt{\tiny{[P1$\leftrightarrow$P]}}
                \texttt{{[}DRAFT{]}} \\
\texttt{No.} \\
\\ 
\texttt{To answer this question, we need to apply the principles of causal inference, specifically the concept of confounding variables and the potential outcomes framework. } \\
\\ 
\texttt{Let's denote the following variables:} \\
\\ 
\texttt{{-} V: Visited England (binary, 0 = no, 1 = yes)} \\
\texttt{{-} S: Skill (unobserved)} \\
\texttt{{-} O: Occupation (binary, 0 = low, 1 = high)} \\
\texttt{{-} Y: Salary (binary, 0 = low, 1 = high)} \\
\\ 
\texttt{We are given the following information:} \\
\\ 
\texttt{1. P(V = 1) = 0.43 (probability of visiting England)} \\
\texttt{2. P(Y = 1 | V = 0) = 0.21 (probability of high salary for those who have not visited England)} \\
\texttt{3. P(Y = 1 | V = 1) = 0.61 (probability of high salary for those who have visited England)} \\
\\ 
\texttt{We want to determine if high salary (Y = 1) is less likely than low salary (Y = 0) overall, i.e., P(Y = 0) > P(Y = 1).} \\
\\ 
\texttt{To do this, we need to consider the potential outcomes framework. Let's denote the potential outcomes as Y0 and Y1, where Y0 is the salary if the individual had not visited England, and Y1 is the salary if the individual had visited England.} \\
\\ 
\texttt{We can use the law of total probability to express the overall probability of high salary as:} \\
\\ 
\texttt{P(Y = 1) = P(V = 0)P(Y1 | V = 0) + P(V = 1)P(Y1 | V = 1)} \\
\\ 
\texttt{Substituting the given values, we get:} \\
\\ 
\texttt{P(Y = 1) = (1 {-} 0.43)(0.21) + 0.43(0.61)} \\
\texttt{= 0.21 + 0.2623} \\
\texttt{= 0.4723} \\
\\ 
\texttt{Now, let's consider the probability of low salary:} \\
\\ 
\texttt{P(Y = 0) = P(V = 0)P(Y0 | V = 0) + P(V = 1)P(Y0 | V = 1)} \\
\\ 
\texttt{Since we don't have direct information about P(Y0 | V = 1), we need to use the fact that P(Y = 1 | V = 1) = 0.61 and P(Y = 0 | V = 1) = 1 {-} P(Y = 1 | V = 1) = 1 {-} 0.61 = 0.39.} \\
\\ 
\texttt{We can also use the fact that P(Y = 1 | V = 0) = 0.21 and P(Y = 0 | V = 0) = 1 {-} P(Y = 1 | V = 0) = 1 {-} 0.21 = 0.79.} \\
\\ 
\texttt{Now, we can use the law of total} \\
            }
        }
    }
    & \\ \\

    \theutterance \stepcounter{utterance}  
    & & \multicolumn{4}{p{0.6\linewidth}}{
        \cellcolor[rgb]{0.95,0.95,0.95}{
            \makecell[{{p{\linewidth}}}]{
                \texttt{\tiny{[P1$\leftrightarrow$P]}}
                \texttt{{[}SELF{-}FEEDBACK{]}} \\
\texttt{FORMAT: ok} \\
\texttt{RULES: ok} \\
\texttt{STRATEGY: problem} \\
\texttt{REASON: The answer does not directly address the question of whether high salary is less likely than low salary overall, but rather calculates the overall probability of high salary, which is not the same thing.} \\
            }
        }
    }
    & \\ \\

    \theutterance \stepcounter{utterance}  
    & \multicolumn{4}{p{0.6\linewidth}}{
        \cellcolor[rgb]{0.561,0.737,0.561}{
            \makecell[{{p{\linewidth}}}]{
                \texttt{\tiny{[P1$\rangle$GM]}}
                \texttt{No.} \\
\\ 
\texttt{To determine if high salary (Y = 1) is less likely than low salary (Y = 0) overall, we need to consider the potential outcomes framework. Let's denote the potential outcomes as Y0 and Y1, where Y0 is the salary if the individual had not visited England, and Y1 is the salary if the individual had visited England.} \\
\\ 
\texttt{We can use the law of total probability to express the overall probability of high salary as:} \\
\\ 
\texttt{P(Y = 1) = P(V = 0)P(Y1 | V = 0) + P(V = 1)P(Y1 | V = 1)} \\
\\ 
\texttt{Substituting the given values, we get:} \\
\\ 
\texttt{P(Y = 1) = (1 {-} 0.43)(0.21) + 0.43(0.61)} \\
\texttt{= 0.21 + 0.2623} \\
\texttt{= 0.4723} \\
\\ 
\texttt{Now, let's consider the probability of low salary:} \\
\\ 
\texttt{P(Y = 0) = P(V = 0)P(Y0 | V = 0) + P(V = 1)P(Y0 | V = 1)} \\
\\ 
\texttt{Since we don't have direct information about P(Y0 | V = 1), we need to use the fact that P(Y = 1 | V = 1) = 0.61 and P(Y = 0 | V = 1) = 1 {-} P(Y = 1 | V = 1) = 1 {-} 0.61 = 0.39.} \\
\\ 
\texttt{We can also use the fact that P(Y = 1 | V = 0) = 0.21 and P(Y = 0 | V = 0) = 1 {-} P(Y = 1 | V = 0) = 1 {-} 0.21 = 0.79.} \\
\\ 
\texttt{Now, we can use the law of total probability to express P(Y0 | V = 1) in terms of P(Y0 | V = 0) and P(Y1 | V = 0):} \\
\\ 
\texttt{P(Y0 | V = 1) = P(Y0 | V = 0) {-} P(Y1 | V = 0) + P(Y1 | V = 1)} \\
\\ 
\texttt{Substituting the given values, we get:} \\
\\ 
\texttt{P(Y0 | V = 1) = 0.79 {-} 0.21 + 0.61} \\
\texttt{= 1.19} \\
\\ 
\texttt{However, this is not possible since P(Y0 | V = 1) must be between 0 and 1. This indicates that there is a problem with our previous assumption that P(Y0 | V = 1) = 1 {-} P(Y = 1 | V = 1).} \\
\\ 
\texttt{To resolve this issue, we can use the fact that P(Y = 1 | V = 0) = 0.21 and P(Y =} \\
            }
        }
    }
    & & \\ \\

    \theutterance \stepcounter{utterance}  
    & & & \multicolumn{2}{p{0.3\linewidth}}{
        \cellcolor[rgb]{0.9,0.9,0.9}{
            \makecell[{{p{\linewidth}}}]{
                \texttt{\tiny{[GM$|$GM]}}
                \texttt{No} \\
            }
        }
    }
    & & \\ \\

    \theutterance \stepcounter{utterance}  
    & & & \multicolumn{2}{p{0.3\linewidth}}{
        \cellcolor[rgb]{0.9,0.9,0.9}{
            \makecell[{{p{\linewidth}}}]{
                \texttt{\tiny{[GM$|$GM]}}
                \texttt{yes} \\
            }
        }
    }
    & & \\ \\

    \theutterance \stepcounter{utterance}  
    & & & \multicolumn{2}{p{0.3\linewidth}}{
        \cellcolor[rgb]{0.9,0.9,0.9}{
            \makecell[{{p{\linewidth}}}]{
                \texttt{\tiny{[GM$|$GM]}}
                \texttt{game\_result = LOSE} \\
            }
        }
    }
    & & \\ \\

\end{supertabular}
}

\end{document}
